% !TeX spellcheck = en_US
\documentclass[french]{yLectureNote}

\title{Algèbre linéaire 2}
\subtitle{Langage mathématique}
\author{Paulhenry Saux}
\date{\today}
\yLanguage{Français}

\professor{C.Dartyge}

\usepackage{graphicx}%----pour mettre des images
\usepackage[utf8]{inputenc}%---encodage
\usepackage{geometry}%---pour modifier les tailles et mettre a4paper
%\usepackage{awesomebox}%---pour les boites d'exercices, de pbq et de croquis ---d\'esactiv\'e pour les TP de PC
\usepackage{tikz}%---pour deiffner + d\'ependance de chemfig
\usepackage{tkz-tab}
\usepackage{chemfig}%---pour deiffner formules chimiques
\usepackage{chemformula}%---pour les formules chimiques en \'equation : \ch{...}
\usepackage{tabularx}%---pour dimensionner automatiquement les tableaux avec variable X
\usepackage{awesomebox}%---Pour les boites info, danger et autres
\usepackage{menukeys}%---Pour deiffner les touches de Calculatrice
\usepackage{fancyhdr}%---pour les en-t\^ete personnalis\'ees
\usepackage{blindtext}%---pour les liens
\usepackage{hyperref}%---pour les liens (\`a mettre en dernier)
\usepackage{caption}%---pour la francisation de la l\'egende table vers Tableau
\usepackage{pifont}
\usepackage{array}%---pour les tableaux
\usepackage{lipsum}
\usepackage{yFlatTable}
\usepackage{multicol}
\newcommand{\Lim}[1]{\lim\limits_{\substack{#1}}\:}
\renewcommand{\vec}{\overrightarrow}
\newcommand{\bz}{\overline{z}}
\DeclareMathOperator{\card}{card}
\renewcommand{\vec}{\overrightarrow}
\newcommand{\N}[0]{\mathbb{N}}
\newcommand{\R}[0]{\mathbb{R}}
\newcommand{\C}[0]{\mathbb{C}}
\newcommand{\dd}[0]{\mathrm{d}}
\begin{document}
\setcounter{chapter}{2}
\chapter{Matrices et Applications linéaires}
\section{Introduction}
\begin{definition}[Mqtrice associées aux applications linéaires]
Soit \(\{e_1,e_2,e_n\}\) une base de E et \(\{\epsilon_1,\dots,\epsilon_p\}\). On écrit \(f(e_1) = a_{11}\epsilon_1+\dots+a_{1p}\epsilon_p\) jusqu'à \(f(e_n) = a_{1n}\epsilon_1+\dots+a_{pn}\epsilon_p\). La mtrice de f dans les 2 bases, notée \(M(f)_{e_i,\epsilon_j}\), la matrice \[
\begin{bmatrix}
a_{11}&p_{12}&\dots& a_{1n}\\
\vdots & &\vdots
a_{p1}&p_{p2}&\dots& a_{pn}\\
\end{bmatrix}\]
\end{definition}
\begin{proposition}[Espace vectoriel de matrice d'application]
Les matrices d'applications sont des isomophismes d'espaces vectoriels et vérifient donc \(M(f+g) = M(f)+M(g), M(\lambda f) = \lambda M(f)\)
\end{proposition}
\begin{proposition}[Produit de 2 MAP]
On a \(M(f\circ g) = M(f)\cdot M(g)\)
\end{proposition}
\begin{proposition}[Propriété de multiplication des matrices]
\(A(BC) = (AB)C, A(B+C) = AB+AC, (A+D)B = AB+DB\) mais \(BA\neq AB\)
\end{proposition}
\section{Matrice d'un vecteur}
\begin{definition}[Matrice d'un vecteur]
Soit \(\{e_1,\dots, e_n\}\) une base de E. La matrice de \(x = x_1e_1+\dots +x_ne_n\) dans cette base est \(M(x)_{e_i}\begin{bmatrix}
x_1\\
\vdots\\
x_n\\
\end{bmatrix}\)
\end{definition}
\begin{definition}[Application de la composition de fonctions à un vecteur]
\(M(f(x)) = M(f)\cdot M(x)\)
\end{definition}

\begin{proposition}[Application linéaire bijective et application réciproque]
Une application linéaire est bijective si sa matrice est inversible. On a alors \(M(f^{-1}) = M(f)^{-1}\)
\end{proposition}
\section{Changement de base}
\subsection{Matrice de passage}
\begin{definition}[Matrice de passage]
Soit 2 bases de E : \(\{e_1,\dots, e_n\}, \{r_1,\dots, r_n\}\). On écrit \(r_1 = p_{11}e_1+\dots p_{n1}e_n, r_n = p{n1}e_1+\dots+p_{nn}e_n\) donc la matrice de changement de base s'écrit :
\[\begin{bmatrix}
p_{11}&p_{12}&\dots& p_{1n}\\
\vdots & &\vdots
p_{p1}&p_{p2}&\dots& p_{pn}\\
\end{bmatrix}\]
\end{definition}
\begin{proposition}[Propriétés des matrices de passage]
Inversibilité : \(P_{e\to v} = P_{v\to e}^{-1}\)

Transitivité : \(P_{e\to v}P_{v\to w} = P_{e\to w}\)
\end{proposition}
\begin{proposition}[Changement de base sur un vecteur]
\[X_v = P_{e\to v}^{-1}X_e\]
\end{proposition}
\begin{proposition}[Changement de base sur une application linéaire]
\[M(f)_{e_1, e_2} = Q_{e_3\to e_2}^{-1}M(f)_{e_3,e_4}P_{e_4, e_1}\]
\end{proposition}
\section{Rang d'une application linéaire}
\begin{proposition}[Rang d'une application et rang de la matrice associée]
Le rang d'une application linéaire est la dimension de son image. On a \(rg(f) = Rg(M(f))\).
\end{proposition}
\begin{proposition}[Rang de la transposée]
\[rg(A) = rg(^tA)\]
\end{proposition}
\subsection{Matrice semblable}
\begin{definition}[Matrice semblable]
2 matrices A et A' sont semblable s'il existe une matrice inversible P telle que \(A' = P^{-1}AP\). On en déduit que deux matrices semblables représentent le même endomorphisme en des bases différentes.
\end{definition}
\warningInfo{Matrice semblable vérification}{Une condition nécessaire est que le déterminant doit \^etre le m\^eme.}

%TODO Méthode pour montrer que 2 bases sont duales.
\subsection{Espace dual}
\begin{definition}[Espace dual]
L'ensemble des formes linéaires \(L_k(E, K)\) est noté E* et est dit espace dual de E. Il a la m\^eme dimension que E.
\end{definition}
%TODO Exe paage 84 à faire. pour trouver base d'un espace dual. + Annulateur
\subsubsection{Trouver une base duale}
\begin{proposition}[Noyau d'une application linéaire dans un espace dual]
\(dim(ker(\omega)) = n-1\) avec n la dimension de E et \(_omega \in E*\).
\end{proposition}
\section{En pratique}
\begin{itemize}
 \item \(Mat(f)_{e_1,e} = Mat(Id)_{e,e_1}Mat(f)_{e,e}Mat(Id)_{e_1,e}\)
 \item \(M_{AB} = M_{\cdot B}\dots M_{A\cdot}\)
 \item Matrice de passage : La matrice de passage de \(e\) à \(e_1\) est la matrice Identité de \(e_1\) à \(e\) : \(Mat(Id)_{e_1,e}\) : on écrit en colonne en fonction de e les vecteurs de \(e_1\)
 \item Matrice de Passage de E à B : On écrit les vecteur colonnes de B en fonction de E : \(M(I_d)_{B,E} = M(I_d)_{E\to B}\)
 \item Donc \(M(f)_{BB} = M(I_d)_{EB}M(f)_{EE}M(I_d)_{BE}\)
\end{itemize}

\end{document}


